\chapter*{Preface}
\addcontentsline{toc}{chapter}{Preface}

Book 4 closed by showing that even total semantic heat death does not extinguish the invariant $\chi_{\mathrm{RSVP}}$, and that a fluctuation reintroducing non-trivial cohomology is always, in principle, available. The present book asks what compels such a fluctuation to occur at all, rather than merely remaining possible. The answer developed here is that the plenum is not neutral with respect to its own persistence: built into the field algebra already established across this series is a compensatory structure, the antifield, whose presence ensures that curvature descent does not merely happen to preserve coherence but is constrained to do so by the same master equation that gives the BV formalism its consistency. This compensatory structure is what the present book calls desire, and the task ahead is to show that the word is not a metaphor imported from ordinary usage but a precise name for an algebraic necessity already implicit in Books 1 through 4.

\chapter*{Diagrammatic Overview}
\addcontentsline{toc}{chapter}{Diagrammatic Overview}

The RSVP triple $\RSVPTriple$, treated throughout the preceding books as the plenum's primitive content, is doubled in this book by the introduction of a dual antifield triple $(\PhiField^{*}, \vField^{*}, \SEntropy^{*})$. Where the original triple describes what the plenum is, the antifield triple describes what the plenum requires of itself in order to remain consistent: each antifield is the algebraic partner that must exist, and must satisfy the master equation developed in Part II, if the corresponding field's curvature descent is not to generate contradiction. The chapters below develop this doubling formally, but the diagram to keep in mind throughout is simple: every field has a shadow, and the shadow is not decorative but load-bearing.

